\subsubsection{The auxiliary character and marked Frobenius}
\label{par:dyadic-residual}

The local chart requires an ordering of the residual characters at several marked places at once. Choosing the auxiliary character first supplies Frobenius and inertia elements that distinguish those characters.

\begin{lemma}[The auxiliary dyadic character and marked elements]\label{lem:marked-dyadic-character}
For the even-order or rank-four construction with
\(n=n_{\rm II}\) or \(n=n_{\rm III}\) of \eqref{eq:ordinary-dyadic-exponents},
respectively, the character \(\eta=\eta_0\eta_2\)
can be chosen unramified at each \(v\in T\). For the indicated arithmetic Frobenius lift \(f_v\) the simultaneous separation is 
\[
 \overline{\eta\Psi^i}(f_v)\ne1\qquad(v\in T,\ i\in\mathbb Z).
\]

The choices of \(g_0\), the prime-to-\(\ell\) Frobenius component
and \(s_0\) below have the residual and cyclotomic properties
required in the graded chart.
\end{lemma}

\begin{proof}
Use the corresponding \(n\) in \eqref{eq:ordinary-dyadic-exponents}.
In the coefficient ring containing the \(2^n\)-th roots of unity take
\(\eta=\eta_0\eta_2\), where \(\eta_2\) is primitive and odd of conductor
\(2^n\). For \(n\ge3\) the unit group is
\(\langle-1\rangle\times\langle5\rangle\), with generator orders
\(2\) and \(2^{n-2}\). Choose
\[
                \eta_2(-1)=-1,\qquad
                \eta_2(5)=\zeta_{2^{n-2}} .
\]
The kernel of the last reduction modulo \(2^{n-1}\)
has generator \(5^{2^{n-3}}\equiv1+2^{n-1}\pmod{2^n}\).
Its character value is
\[
       \eta_2(5^{2^{n-3}})=\zeta_{2^{n-2}}^{\,2^{n-3}}=-1,
\]
which proves conductor \(2^n\).
In the even-order case take \(\eta_0=1\).

In the rank-four case the nonprincipal character \(\xi\) of \(U/H\) has odd order dividing \(p_1-1\) and prime to \(\ell\).
Take \(\eta_0\) of that order modulo \(p_1\), possible by cyclicity
and the presence of the values of \(\xi\) in the coefficient ring. It is even.
In both cases \(\eta\) has order prime to \(\ell\).
For \(v\ne\ell\) the \(\ell\)-part of the inertia character factors through \(\mathbf F_v^\times\). Thus
\[
 v\in S_b\quad\Longrightarrow\quad \ell\mid |\mathbf F_v^\times|=v-1 .
\] In the rank-four case this set is empty.
Thus \(\eta\) is unramified at every \(v\in T\). Choose arithmetic
Frobenius lifts there, at \(\ell\) choosing one with \(\epsilon_\ell=1\)
by total ramification of the cyclotomic tower. The residual cyclotomic
value of each lift is one.

For \(v\in T\) and \(i\in\Z\) we claim
\(\overline{\eta\Psi^i}(\operatorname{Frob}_v)\ne1\).
In the even-order case an exponent \(2^a u\) killing \(\xi,\eta_0\) may be taken
with
\[
 a=c_2(\xi)+1+v_2(o)
\]
and \(u\) odd, by factoring the host
character into its dyadic, odd-conductor prime-to-\(\ell\), and
\(\ell\)-primary parts. In the rank-four case take \(a=0,\ u=\operatorname{ord}(\xi)\).
If the asserted value were one, raising it to this exponent would give
\(\eta_2(v)^{2^a u}=1\), since reduction on roots of 2-power order is
injective. Put \(s=v_2(v^2-1)\). As \(u\) is odd and \(v^2\equiv1\pmod8\), the odd geometric sum gives
\[
 v_2(v^{2u}-1)=s.
\]
Successive squaring now uses
\[
 v_2(x^2-1)=v_2(x-1)+v_2(x+1)=v_2(x-1)+1
 \qquad(x\equiv1\bmod 4).
\]

In particular \(v_2(v^{2^{a+1}u}-1)=s+a\). If \(s+a<n\), the repeated square gives
\[
    (v^2)^{u2^a\cdot2^{n-s-a-1}}\equiv1+2^{n-1}\pmod{2^n}.
\]
This follows because \(v^2\equiv1\bmod8\) and the valuation
increases by one at each squaring. The character value there
is \(-1\), contrary to the assumed vanishing. Thus vanishing on \((v^2)^{u2^a}\) would force
\[
         2^n\le2^{s+a}\le(v^2-1)2^a<v^22^a .
\]
In the even-order case, \(v\le\ell N\) and \(2^{v_2(o)}\le o\)
give \(v^22^a\le 2^{c_2(\xi)+1}o\ell^2N^2\le 2^n\).
In the rank-four case \(v=\ell\), \(a=0\), and
\(\ell^2\le2^n\). Both contradict the strict inequality.

In both cases the prime-to-\(\ell\) factor of the host character
after removing the factor at \(\ell\) is nontrivial. It cannot be
cancelled in \(\bar\Psi\) by a Teichm\"uller power: restriction
to inertia at \(\ell\) would force that power to be trivial first.
Thus \(\bar\Psi\) is nontrivial, and so is \(\bar\alpha\) by the
marked value just proved. If neither kernel contains the other, take
\(g\in\ker\bar\alpha\setminus\ker\bar\Psi\) and
\(h\in\ker\bar\Psi\setminus\ker\bar\alpha\). Then
\[
    \bar\alpha(gh)=\bar\alpha(h)\ne1,\qquad
    \bar\Psi(gh)=\bar\Psi(g)\ne1 .
\]

In the containment case choose an
element outside the larger proper kernel. This gives the required
\(g_0\).
In the closure of the cyclic group generated by the Frobenius at
\(\ell\), take the prime-to-\(\ell\) component. It is a limit of
powers (use exponents \(\ell^{a!}\) as \(a\to\infty\)) and retains
the residual character values. Choose an inertia element with
cyclotomic value of infinite order and raise it to an integer
divisible by the orders of \(\xi,\omega\). The resulting \(s_0\)
kills both finite factors and still has infinite cyclotomic order.

\end{proof}

\subsubsection{Vanishing for the six odd roots}

The residual ordering fixes the action on each off-diagonal root space. Before counting parameters in the chart, we must decide which root cocycles survive the marked local conditions modulo coboundaries.

\begin{proposition}[Odd-root cocycles]\label{prop:six-odd-roots}
For the six residual characters
\(\bar\alpha,\bar\beta,\bar\vartheta\) and their inverses constructed
in Lemma~\ref{lem:marked-dyadic-character}, every cocycle \(z\) for such a character \(\theta\) satisfying the marked local conditions of the graded chart is a coboundary: 
\[
 \exists b\quad z(g)=(\theta(g)-1)b\qquad\text{for every }g.
\]

\end{proposition}

\begin{proof}
\proofstep{The six primitive odd characters.}

The characters in question are \(\bar\alpha,\bar\beta,\bar\vartheta\)
and their inverses. Each is the reduction of \(\varphi\rho\) where
\(\varphi\) is primitive of odd conductor \(f\) and \(\rho\) primitive
of conductor \(2^n\). Indeed \(\bar\alpha=\bar\eta\), and the other
two are products of \(\bar\Psi\) and \(\bar\eta^{\pm1}\). The dyadic factor of \(\xi\) has conductor
\(2^{c_2(\xi)}<2^n\). On the last reduction kernel its
restriction is trivial. Consequently the dyadic products preserve
\[
          \rho(1+2^{n-1})=-1 .
\] The odd factors
are formed from the prime-to-\(\ell\) odd-conductor host part, a
Teichm\"uller power at \(\ell\), and \(\eta_0^{\pm1}\).

Let \(d\)
be the multiplicative order used in defining \(n_0\) below, with
\(r=v_2(\ell^d-1)=n_0\ge2\). These factors have order dividing
\(\ell^d-1\), also for \(\eta_0\) by its prescribed order. Their
primitive product has conductor dividing \(b\ell\) times the
conductor of the prime-to-\(\ell\) odd-conductor host part. Hence
\(2^n\ge f4^r\) by the first term in the definition of \(n\).
The products are odd by \(\xi(-1)=(-1)^j\) and the oddness of \(\eta\).

Set \(h=2^r,\ L_0=2^{n-r},\ Q_0=f2^n\). The conductor inequality yields
\[
 L_0=2^n/h\ge fh,\qquad v_2(L_0^2)=2(n-r)\ge n .
\]
In particular \((1+L_0)^i\equiv1+iL_0\pmod{2^n}\). The root \(\zeta=\rho(1+L_0)\)
has order \(h\), since \(1+L_0\) generates the kernel of reduction
to \(L_0\) and \(\rho\) is primitive. Take one odd representative
\(1\le u<L_0\) of each unit class modulo \(L_0\) up to sign,
and put \(z_u=\zeta^{u^{-1}\bmod h}\). Pair the summands at \(a\)
and \(Q_0-a\) in
\[
 B_{1,\varphi\rho}=Q_0^{-1}\sum_a a\varphi(a)\rho(a).
\]

The weight is now \(2a-Q_0\). Write \(a=u+L_0(i+fs)\) with
\(0\le i<f,\ 0\le s<h\). The dyadic factor is \(\rho(u)z_u^{i+fs}\), and
\(\omega=z_u^f\) has order \(h\) since \(f\) is odd.
The constant and linear parts of the weight use the respective sums
\[
       \sum_{s=0}^{h-1}\omega^s=0,\qquad
       \sum_{s=0}^{h-1}s\omega^s
          =\left.X\frac{d}{dX}\frac{X^h-1}{X-1}\right|_{X=\omega}
          =\frac{h}{\omega-1}.
\] Now
\(2L_0 f h/Q_0=2\), and summing in \(i\) gives
\[
 B_{1,\varphi\rho}=2\sum_u \rho(u)\,
       \frac{\sum_{i=0}^{f-1}\varphi(u+L_0i)z_u^i}{z_u^f-1}.
\]
Every denominator is a unit at \(\ell\), even when \(\ell\mid f\).
After reduction the fractions are in \(k=\mathbf F_{\ell^d}\).
The element \(\bar\rho(5)\) has order \(2^{n-2}\). From
\(2^n\ge f4^r\) and \(r\ge2\) we have \(n-2\ge r\). The least
\(m>0\) with \(2^{n-2}\mid(\ell^d)^m-1\) is
\[
         D=2^{n-r-2},
        \qquad v_2((\ell^d)^m-1)=r+v_2(m).
\]

Thus the degree over \(k\) is \(D\).
The powers \(5^j\) for \(0\le j<D\) enumerate the classes modulo
\(L_0\) up to sign. The values \(\bar\rho(u)\) differ from the
corresponding \(\bar\rho(5)^j\) by elements of \(k^\times\).
They are linearly independent. Vanishing modulo \(\ell\) would
force each numerator to vanish for every choice of representatives.
Fix an odd class of \(u\bmod h\), so \(z=\bar z_u\) is fixed.
Since \(L_0\ge fh\), such \(u\) in the allowed interval cover
every class modulo \(f\). Thus
\(F(x)=\sum_{i=0}^{f-1}\bar\varphi(x+L_0i)z^i\) would vanish
identically. Shifting the sum gives
\[
        zF(x+L_0)-F(x)=(z^f-1)\bar\varphi(x),
\]
which is nonzero at \(x=1\). This proves the unit assertion,
also with both characters inverted.
\hfill\mbox{\cite{washington-nonp}}\label{cor:ordinary-residual-characters-1}

\proofstep{Use the integral Stickelberger annihilator.}

The marked cocycles have been reduced to character components of a class-group calculation. An integral annihilator is needed here because rational vanishing alone leaves possible residual torsion.


Put \(\chi=\varphi\rho\) and \(\gamma=\bar\chi\). A residual
cocycle satisfying the chart conditions is a homomorphism on
\(\ker\gamma\). It is zero there on inertia. At \(T\) the cocycle
is locally a coboundary. Elsewhere it kills the inertia kernel
of the diagonal residual representation, whose index is prime
to \(\ell\). On \(G_{\Q(\mu_{Q_0})}\) it therefore factors through
the ideal class group by class field theory, equivariantly for
\(\gamma\). Choose \(a\equiv1\bmod\ell f\) and
\(a\equiv1+2^{n-1}\bmod 2^n\). Then \(\bar\chi(a)=-1\).

For
\[
 \Theta=Q_0^{-1}\sum_b b\sigma_b^{-1},
\]
with \(b\)
running through least positive unit representatives, the integral
Stickelberger element \((1-a\sigma_a^{-1})\Theta\) annihilates
the cyclotomic class group. The two contributions to the coefficient at \(\sigma_b^{-1}\) are
\[
 \frac{b}{Q_0}-\frac{a[b/a]_{Q_0}}{Q_0}\in\Z ,
 \qquad a[b/a]_{Q_0}\equiv b\pmod {Q_0}.
\]
Here the brackets select the least positive residue. Evaluation gives
\[
         (1-a\chi(a)^{-1})B_{1,\chi^{-1}}.
\]
Since \(a\equiv1\pmod\ell\) and \(\bar\chi(a)=-1\), the first
factor reduces to \(2\). The second is a unit by the preceding
calculation.
The homomorphism on the class group is consequently zero. The
cocycle now factors through an abelian group. For an element \(s\) with \(\gamma(s)\ne1\), compare the
cocycle equation at \(gs=sg\):
\[
      (\gamma(s)-1)z(g)=(\gamma(g)-1)z(s).
\]
This gives the coboundary in each of the six degrees.
\hfill\mbox{\cite{washington-cyclo}}\label{cor:ordinary-residual-characters-2}
\end{proof}
